\subsection{Greedy Best First Search (GBFS)}

\textbf{Ý tưởng:} GBFS là một chiến lược tìm kiếm thuộc nhóm \textbf{Informed-Search}, hướng tới mục tiêu bằng cách luôn chọn mở rộng nút \textbf{có vẻ gần đích nhất} theo một hàm ước lượng heuristic $h(n)$. GBFS \textbf{chỉ quan tâm đến chi phí ước tính còn lại} đến mục tiêu. \cite{russell2021}

\textbf{Cấu trúc dữ liệu:} \texttt{priority queue} (min-heap), ánh xạ \texttt{visited[node] = parent}, hàm heuristic \texttt{h(n)}.

\textbf{Cơ chế thực hiện:}
\begin{itemize}
    \item GBFS mở rộng nút $n$ có giá trị \textbf{heuristic nhỏ nhất} trong hàng đợi ưu tiên (priority queue).  
    \item Hàm đánh giá sử dụng là:
    \[
        f(n) = h(n)
    \]
    trong đó:
    \begin{itemize}
        \item $h(n)$ là \textbf{ước tính chi phí nhỏ nhất} từ nút $n$ đến mục tiêu. Tại đây, $h(n)$ trả về giá trị trọng số giữa đỉnh đang xét tới láng giềng của nó, trọng số càng nhỏ, $h(n)$ càng nhỏ, láng giềng có $h(n)$ nhỏ nhất sẽ được thêm vào frontier:
        \[
        h(n) = weight[n][neighbor] = matrix[n][neighbor]
        \]
        \item Không xét đến chi phí thực tế $g(n)$ (chi phí từ gốc đến $n$).
    \end{itemize}
    \item Ở mỗi bước, GBFS chọn nút có $h(n)$ nhỏ nhất để mở rộng, sinh ra các nút con và tính giá trị heuristic mới cho chúng.
    \item Kiểm tra mục tiêu được thực hiện khi một nút được chọn để mở rộng.
    \item Quá trình kết thúc khi tìm thấy trạng thái mục tiêu hoặc hàng đợi rỗng (không còn trạng thái có thể mở rộng).
\end{itemize}

\textbf{Ghi chú triển khai:}
\begin{itemize}
    \item GBFS sử dụng \textbf{priority queue (min-heap)} để lưu các nút theo giá trị $h(n)$ tăng dần.
    \item Nếu có nhiều nút cùng giá trị $h(n)$, thuật toán có thể áp dụng \textbf{tie-break rule}.
    \item Hiệu quả của GBFS phụ thuộc mạnh vào chất lượng của \textbf{hàm heuristic}: 
    \begin{itemize}
        \item Nếu $h(n)$ gần với chi phí thật, GBFS tìm đích rất nhanh.
        \item Nếu $h(n)$ sai lệch, thuật toán có thể đi lạc hướng hoặc quay lui nhiều lần.
    \end{itemize}
\end{itemize}

\textbf{Đánh giá Hiệu suất:}
\begin{itemize}
    \item \textbf{Tính đầy đủ (Completeness):} Có, nếu không gian trạng thái hữu hạn (finite graph search).
    \item \textbf{Tính tối ưu (Optimality):} Không, vì GBFS có thể bỏ qua lời giải tối ưu khi heuristic đánh giá sai hướng.
    \item \textbf{Độ phức tạp thời gian (Time Complexity):} $O(|V|)$, 
    trong đó $|V|$ là số đỉnh của đồ thị trong trường hợp xấu nhất. Tuy nhiên, với một hàm heuristic tốt, 
    số lượng nút mở rộng có thể giảm đáng kể; trong một số bài toán, độ phức tạp thực tế đạt khoảng $O(bm)$, với $b$ là hệ số phân nhánh và $m$ là độ sâu của lời giải.
    \item \textbf{Độ phức tạp không gian (Space Complexity):}  $O(|V|)$, vì GBFS cần lưu tất cả các nút đã sinh ra trong hàng đợi ưu tiên (frontier) để tránh lặp lại.
    \item \textbf{Nhận xét:}
    \begin{itemize}
        \item Do không xét $g(n)$, GBFS có thể chọn đường “ngắn trước mắt” nhưng “tốn kém tổng thể”.
        \item GBFS thường chạy nhanh hơn BFS và UCS trong các bài toán có không gian lớn, nhưng không đảm bảo tối ưu — do đó nó thích hợp cho các bài toán cần \textbf{tìm nhanh một lời giải chấp nhận được} thay vì lời giải tối ưu.
    \end{itemize}
\end{itemize}